\documentclass{article}
\usepackage{/globals/math}
\usepackage{/globals/code}
\usepackage{/globals/tables}
\usepackage{/globals/anki}
\ankishow
\ankiprefix{datatypes}
\ankitopic{Datentypen}
\begin{document}
\section{Datentypen} 
Welche Werte eine Variable annehmen kann, was mit diesen gemacht werden kann und was diese repräsentieren wird durch den \emph{Datentypen} angegeben.

\subsection{Ganze Zahlen}
Ganze Zahlen können in vier unterschiedlichen Datentypen gespeichert werden; \ilsourcecode{byte}, \ilsourcecode{short}, \ilsourcecode{int} und \ilsourcecode{long}. Diese unterscheiden sich darin, wie viele Bytes sie an Arbeitsspeicher einnehmen und dementsprechend auch wie groß die darin gespeicherten Zahlen sein können.
\begin{ztable}{lll}{Datentype & Größe & Wertebereich}
 byte & 1 Byte (8 Bit) & $-2^7$ bis $2^7-1$ \\
 short & 2 Byte (16 Bit) & $-2^{15}$ bis $2^{15}-1$ \\
 int & 4 Byte (32 Bit) & $-2^{31}$ bis $2^{31}-1$ \\
 long & 8 Byte (64 Bit) & $-2^{63}$ bis $2^{63}-1$ \\
\end{ztable}
Mit unterstrichen können visuelle Seperatoren dargestellt werden, so wie Punkte im Deutschen oder Kommas im Englischen genutzt werden. Zudem kann mit einem \ilsourcecode{e} die wissenschaftliche Notation genutzt werden. Dieses beide gilt auch für Gleitkommezahlen.
\begin{slsourcecode}{java}
int a = 1_000_000;
int a = 1e6; // 1 * 10^6
\end{slsourcecode}
Standardmäßig werden \ilsourcecode{int}s genutzt. Dass eine Zahl ein \ilsourcecode{long} ist kann mit dem \ilsourcecode{L} suffix angemerkt werden.
\begin{slsourcecode}{java}
int a = 2_147_483_647;
long b = 2_147_483_648L;
\end{slsourcecode}
Wird du der größtmöglich speicherbaren Zahl eins addiert, so kommt es zu einem Überlauf; sie wird zur kleinstmöglichen Zahl. Gegenteilig ist die kleinstmögliche Zahl minus eins die größtmöglich Zahl.
\begin{slsourcecode}{java}
int a = 2_147_483_647;
System.out.println(a + 1); // -2147483648
\end{slsourcecode}

\ankinote{sizes-bits}{Zahlen Größen Bits}{byte = 1 Byte, 8 Bit, danach doppelnd}
\ankinote{sizes-vals}{Zahlen Wertebereiche}{byte = \(2^8\), danach doppelnd}
\subsection{Dezimalzahlen}
Dezimalzahlen werden als \ilsourcecode{float}s und \ilsourcecode{double}s dargestellt, jeweils 4 Bytes und 8 Bytes in der größe.

Sie werden mit einem Punkt als Komma aufgeschrieben, so wie Dezimalzahlen auch im Englischen geschrieben werden. Mit einem \ilsourcecode{f} kann gekennzeichnet werden dass eine Zahl ein \ilsourcecode{float} ist.

Gleitkommazahlen funktionieren indem ein Bit das Vorzeichen und jeweils ein paar Bits für die Mantisse und den Exponenten genutzt werden. Dann ist
\[
 \text{Wert} = \text{Vorzeichen} * \text{Mantisse} * 2^{\text{Exponent}}
\]
Hier ist das Vorzeichen $1$ oder $-1$ und die Mantisse zwischen $1$ und $2$. So können Dezimalzahlen unterschiedlichster Größenordnungen.

\ankinote{float}{Float Bestandteile}{Vorzeichen, Mantisse, Exponent}

\subsection{Wahrheitswerte}
Ein boolean kann entweder Wahr (\ilsourcecode{true}) oder Falsch (\ilsourcecode{false}) speichern.

\subsection{Zeichen}
Ein Charakter, ein \ilsourcecode{char}, kann in zwei Bytes Buchstaben und andere Symbole speichern. Dabei werden einfache Anführungszeichen genutzt.
\begin{slsourcecode}{java}
char c = 'a';
\end{slsourcecode}
Ein Backslash (\textbackslash) zeigt eine Escape-Sequenz an; zusammen mit einem nachfolgenden Zeichen können so Sonderzeichen geschrieben werden.
\begin{ztable}{lll}{Squenz & Sonderzeichen}
 \textbackslash n & Neue Zeile \\
 \textbackslash t & Tabulator \\
 \textbackslash ' & Einfaches Anführungszeichen \\
 \textbackslash " & Doppeltes Anführungszeichen \\
 \textbackslash \textbackslash & Backslash \\
 \textbackslash u<4 Byte Hex Zahl> & Unicode Charakter \\
\end{ztable}

\ankinote{char-u}{Char Unicode}{escape u<4 Byte Hex Zahl>}

\subsection{Text}
Ein \ilsourcecode{String} ist eine Zeichenkette und kann so ganze Text darstellen.
\begin{slsourcecode}{java}
String s = "Hello World!";
\end{slsourcecode}
\end{document}